\section{An evidence-triggered roadmap from RAKL 3.0 to later challengers}
\label{sec:roadmap}

A roadmap can itself become a source of confirmation bias if version numbers are assigned to ideas before evaluation. We therefore distinguish \emph{candidate sequence} from \emph{released method sequence}. A label such as 3.1 is earned only after the corresponding challenger passes its governed process. If the first candidate fails, it remains a rejected branch and a different challenger may eventually become 3.1.

\subsection{3.0.x: measurement and authority hardening without claiming a new research method}

Several immediate changes are instrumentation or implementation hardening rather than intended research-policy innovations. Suitable patch candidates include:

\begin{itemize}[leftmargin=*]
\item read-only state/process metrology and the four-arm attribution data structures introduced on the Paper 5 branch;
\item content-bound replacement for the v3 archive's bare governance Boolean while preserving the existing constitutional requirement that promotion needs governance approval;
\item content resolution of verification/review/proof IDs before authority-bearing use;
\item certificate-backed local-section verification on gluing authority paths;
\item promotion-grade binding of benchmark state/output identities to exact task/evaluator/artifact subjects;
\item public API and CI stabilization.
\end{itemize}

These repairs should improve observability and enforcement of already stated principles. Unless they intentionally change search behavior, they do not by themselves support a new method-generation claim.

\subsection{First Class-B candidate: coverage-bound research memory}

The strongest currently observed cross-problem process defect is the XM002 retrieval miss: a synthesis stated that a Hodge reuse remained absent even though the exact base already contained the corresponding ResearchMemoryReview and reuse artifacts. Framework issue \texttt{RAKL\#119} therefore proposes a content-bound cross-problem coverage receipt.

The candidate mechanism is straightforward. A memory query binds the exact problem/lane universe, index/manifests and structural query coordinates before making completeness-like statements. The primary prediction is lower missed-relevant-memory rate without requiring full-repository loading. Secondary predictions include lower duplicate research and better cross-domain tool reuse. Main risks are context/compute overhead and false confidence in an incomplete index.

The cheapest discriminating benchmark uses frozen repository snapshots with planted relevant and near-miss memories across several lanes, comparing the current query path with the coverage-bound challenger at matched retrieval budget. A successful development result still requires fresh surface-shifted tasks before the candidate can become a new method version.

\subsection{Root-coordinate preservation audit}

Several live case-study episodes exhibit a shared abstraction: an appealing local or surrogate quantity does not preserve the root-critical coordinate. Examples include raw closure mass versus cover complexity, positive norm versus Li-coefficient faithfulness, augmentation order versus complex analytic order, finite-cutoff spectral information versus continuum physical gap, and detector visibility versus actual geometric obstruction vanishing.

If the pattern survives enough independent episodes, a Class-B challenger can require an explicit \texttt{RootCoordinatePreservationAudit} before a surrogate becomes a major route. The intervention asks for the exact map from surrogate to root coordinate, non-preservation examples, cheapest hostile control and remaining gluing obligations. Its primary metric is reduction in root-coordinate false-progress/candidate-generation rate. Its main risk is excessive conservatism that blocks useful exploratory representations before a bridge can be discovered.

\subsection{Gluing-aware search policy}

Another recurring pattern is local mathematical success with an unresolved global interface. Yang--Mills provides an abstract spectral lemma whose target-specific OS/density/continuum binding remains open; Navier--Stokes repeatedly exposes far-field, compactness or limit-passage interfaces after valid local estimates. V3 represents gluing obstructions, but a later challenger could make interface risk an explicit routing coordinate.

The predicted benefit is earlier detection of ``correct but non-closing'' local work and higher residual contraction per cycle. The hostile control is a problem where local decomposition is sufficient and extra gluing machinery only adds overhead. The candidate should fail if it systematically diverts effort from locally solvable tasks.

\subsection{Experience-conditioned routing and strategy motifs}

V3 already contains experience-conditioned operator scoring and motif induction. The correct next action is evaluation, not automatic elaboration. Once enough v3 telemetry accumulates, compare routing with and without experiential priors on repeated-family, transfer and hostile near-miss tasks. A successful result should reduce repeated failures and/or cost while maintaining false-transfer control. If gains disappear under sham memory or fresh transfer, the mechanism should remain local or be revised rather than promoted as a general method improvement.

\subsection{Consolidation policy as an evolvable object}

The v3 fast/slow split intentionally prevents every episode from being immediately rewritten into a lesson. Later data may support different consolidation schedules by failure type, task family or evidence density. This is a natural candidate for method learning because external research has demonstrated that forced continuous consolidation can degrade memory. Any challenger must preserve raw episodes and be compared against episodic-only and current gated-consolidation controls.

\subsection{Portfolio and saturation policy}

The long-horizon autonomous-research literature and our own saturation vector both suggest a failure mode in which an agent repeatedly hill-climbs within one method family after returns flatten. A future challenger could use measured axis-specific saturation and stable residuals to reallocate budget among exploitation, diversification, moonshots and meta-method search. The benchmark must include tasks where persistence within one family is actually correct so that diversification is not rewarded merely for being novel.

\subsection{Meta-method invention}

Only after relevant operator/path/meta-method axes are demonstrably saturated and stable residuals remain should the system search for a genuinely new research operator or ontology. The v3 invention gate operationalizes this principle. Paper 5 treats operator invention as a measured exception rather than the default explanation of progress. This creates an empirical test of the RAKL-triviality hypothesis: how often do verified solutions require new primitive research structure rather than recomposition or transfer?

\subsection{What later versions should publish}

Each released method version should have a compact ``model card'' style evidence packet:

\begin{itemize}[leftmargin=*]
\item exact parent and candidate Git/method identities;
\item source pathologies or positive motifs motivating the change;
\item preregistered mechanism and predicted metrics;
\item matched development results and uncertainty;
\item fresh protected assurance results;
\item regressions and baseline-only wins;
\item resource profile;
\item state/schema migration and rollback status;
\item evaluator/assurance identities and exposure history;
\item final governance/promotion and active-main attestation.
\end{itemize}

Paper 5 can then become a longitudinal scientific record of the sequence
\[
\mathrm{RAKL}_{3.0}\rightarrow \mathrm{RAKL}_{3.1}\rightarrow \mathrm{RAKL}_{3.2}\rightarrow\cdots
\]
without assuming in advance that every arrow points to a better system.